% =============================================================================
%  Front matter: purpose, notation, legend
% =============================================================================
\section*{Purpose}
A one-dimensional boundary-layer scheme is a stack of assumptions, each of which was
reasonable for a flat, homogeneous surface at a grid spacing of ten kilometres. This note
splits that stack into its smallest testable statements (nineteen rows), shows where the
original MYNN~2.5 scheme of \textcite{nakanishi2009} uses each of them, why each fails over
mountains and in the grey zone, and how eleven other schemes handle it: kept, relaxed, or
dropped --- and if dropped, what is assumed instead. It comes in three versions: a one-page
quicklook of what each scheme does (Part~0), an overview with the mechanisms as equations
(Part~I), and a lookup with constants, equation numbers and code references (Part~II). Section~I.3 (added 2026-09-12) marks the rows the PhD concept turns on and judges whether the assumptions the concept names are the right ones and which are missing.
Everything stated about the WRF schemes was read from the local WRF~4.8 sources
(\texttt{branko}, \texttt{goger18WRF}, \texttt{goger19WRF}); everything about the published
schemes from the papers, with equation numbers.

\section*{Notation}
{\small Every symbol used in the tables, grouped by role. Where two schemes use one letter for
different things, both meanings are listed with the context in which they occur. Units are SI;
"--" marks a dimensionless quantity; "code" marks a Fortran variable named as in the source.\par}
\footnotesize
\setlength{\tabcolsep}{4pt}
\renewcommand{\arraystretch}{1.1}
\newcommand{\notgroup}[1]{\multicolumn{3}{@{}l}{\rule{0pt}{3ex}\textbf{#1}} \\ \addlinespace[1pt]}
\begin{xltabular}{\linewidth}{@{}>{\raggedright\arraybackslash}p{3.6cm} >{\raggedright\arraybackslash}X >{\raggedright\arraybackslash}p{2.3cm}@{}}
\toprule
\textbf{Symbol} & \textbf{Meaning} & \textbf{Unit} \\
\midrule
\endfirsthead
\toprule
\textbf{Symbol} (continued) & \textbf{Meaning} & \textbf{Unit} \\
\midrule
\endhead
\bottomrule
\endlastfoot
\notgroup{Indices, subscripts and decorations}
$i, j, k, n$ (on $u$, $U$, $x$, $\tau$, $g$, $f$) & Cartesian component indices, $1, 2, 3$ for $x, y, z$; a repeated index is summed ($P_s = -\avg{u_iu_j}\,\pd{U_i}{x_j}$ runs over nine terms). On its own, $k$ is also the vertical level index of a model column. & -- \\
$U_x, U_y, U_z, V_x, \ldots, W_z$; $\Theta_{v,z}$ & Shorthand for the partial derivative of a mean field: $U_x \equiv \pd{U}{x}$, $\Theta_{v,z} \equiv \pd{\Theta_v}{z}$. & s$^{-1}$; K\,m$^{-1}$ \\
$\partial_t$, $\partial_x$, $\partial_y$, $\partial_z$; $\Delta_x(\cdot)$, $\Delta_z(\cdot)$ & Partial derivatives; finite differences across one grid interval in $x$ or in the vertical, as written in the code. & -- \\
$\avg{\cdot}$ & Grid-cell average: an ensemble average in the 1D closures and the 3D Mellor--Yamada closure, a spatial filter in Ito et al.\ and the LES-type closures (row A2). & -- \\
$\overline{\cdot}$ & Average of a staggered quantity onto the grid point where it is used (e.g.\ $\overline{D_{12}^2}$, $\overline{|z_x|}$: four-point averages); in quoted formulas of Zhang et al.\ and Ito et al.\ the overbar is the resolved (filtered) field. & -- \\
$u', \theta', w'\theta'$ (primed) & Subgrid fluctuation, in the notation of the quoted papers; identical to the unprimed lower-case $u_i$, $\theta$ of the Mellor--Yamada notation used elsewhere. & as the variable \\
subscript $0$: $\avg{w\theta}_0$, $\avg{w\theta_v}_0$, $\avg{wq}_0$, $\theta_0$, $T_0$, $\rho_0$, $\Pr_0$, $l_0$ & Value at the surface (the surface fluxes), or a reference value ($\theta_0 = T_0 = 300$~K, $\rho_0$ the reference density); $\Pr_0$ the surface-layer Prandtl number of YSU; $l_0$ Blackadar's asymptotic length (see below). & as the variable \\
subscript $1$: $u_1$, $v_1$, $U_1$, $V_1$, $z_1$, $\theta_{v,1}$, $\rho_1$, $\Delta z_1$ & Value at, or thickness of, the lowest model level. & as the variable \\
subscript $h$ on a flux: $\avg{w\theta}_h$, $\avg{w\theta_v}_h$, $\avg{wc}_h$ & Entrainment flux at the boundary-layer top $h$ (YSU, Shin--Hong). & as the flux \\
subscripts $h$, $v$ on $K$, $l$, $\ell$, $\Delta$: $K_h$, $K_v$ (rows B5, C2), $l_h$, $\ell_h$, $\ell_v$, $\Delta_h$ & Horizontal and vertical. Exception: in the YSU and Shin--Hong cells $K_m$, $K_h$, $K_c$, $K_q$ follow those schemes' convention and mean the eddy diffusivity for momentum, heat, a generic scalar $c$, and moisture. & -- \\
subscripts $mh$, $mv$, $hh$, $hv$: $K_{mh}$, $K_{mv}$, $K_{hh}$, $K_{hv}$ & WRF's diffusion code: first letter momentum ($m$) or heat ($h$), second letter horizontal ($h$) or vertical ($v$). & -- \\
subscripts $M$, $H$, $q$: $K_M$, $K_H$, $S_M$, $S_H$, $S_q$, $G_M$, $G_H$ & Mellor--Yamada convention: momentum, heat, turbulence kinetic energy. & -- \\
subscripts $S$, $T$, $B$, $F$ on $L$ or $l$ & Surface-layer, turbulent-layer, buoyancy, and free-atmosphere length scales of the MYNN blend. & -- \\
subscripts $up$, $down$, $k$, $\varepsilon$ on $l$ & BouLac: upward and downward parcel displacements, the mixing length, the dissipation length. & -- \\
subscripts $LES$, $meso$ or $M$, $diss$ or $dissip$, $scale$, $master$, $use$, $eff$, $G$ on $l$ or $\ell$ & Length scale in the LES limit, in the mesoscale limit, for dissipation, from WRF's \texttt{calc\_l\_scale}, the fork's master scale, the strain-limited scale actually used, the G19 port's effective and unblended (Goger) horizontal scales. & -- \\
subscripts $NL$, $L$, $u$, $q$, $\theta$, $nl$, $TKE$ on $P$ & Partition (scale-awareness) functions for the nonlocal heat flux, the local heat flux, momentum, moisture, heat, the nonlocal term (Zhang et al.), and TKE. & -- \\
subscript or superscript $*$: $u_*$, $w_*$, $\theta_*$, $q^*$, $L^*$, $\Delta^*$ & Surface-layer and convective scales (friction velocity, convective velocity, temperature scale); in Ito et al.\ a quantity made dimensionless with $w_*$ and $h$. & -- \\
subscripts $c$, $eq$, $2$, $HL88$, $min$, $div$ on $q$ or $\alpha$ & $q_c$: NN09's convective velocity scale; $q_{eq} = q_2 = q_{HL88}$: the level-2 equilibrium value of $q$; $q_{min}$: the numerical floor; $q_{div}$: the Helfand--Labraga ratio $q/q_{eq}$; $\alpha_c$: the resulting disequilibrium factor. & -- \\
subscripts $ts$, $te$: $k_{ts}$, $k_{te}$ & WRF tile indices of the lowest and the highest model level. & -- \\
subscripts $x$, $y$ on $z$ and $\tau$: $z_x$, $z_y$, $\tau_x$, $\tau_y$ & Slope of the model (terrain-following) surface, $(\pd{z}{x})_\eta$; components of the fork's slope factor. & -- \\
subscripts $u$, $v$, $w$ on $\sigma$; $L,u$ and $L,v$ on $T$ & Standard deviations of the velocity components; streamwise and spanwise Lagrangian integral time scales. & m\,s$^{-1}$; s \\
subscripts $mf$, $bl$, $shcu$, $ent$, $sf$, $ns$, $fac$, $lim$, $conv$ & Mass flux ($l_B^{mf}$, $q_{mf}$), boundary layer ($\Psi_{bl}$), shallow cumulus ($\Psi_{shcu}$), entrainment ($f_{ent}$, $c_{ent}$), surface ($c_{sf}$), stable ($c_{ns}$), factor ($\phi_{fac}$, $a_{2fac}$), limit ($\Pr_{lim}$), convective ($v_{conv}$). & -- \\
superscripts $S$, $mf$, $LES$ on $K$ or $l$: $K^S_{mh}$, $l_B^{mf}$, $\ell_v^{LES}$ & The Smagorinsky part of a coefficient, the mass-flux contribution to a length, the LES-limit value. & -- \\
\notgroup{Mean fields and turbulent fluctuations}
$U, V, W$; $U_i$ & Mean (resolved) wind components along $x$, $y$, $z$; index form. & m\,s$^{-1}$ \\
$u, v, w$; $u_i$ & Subgrid velocity fluctuations. & m\,s$^{-1}$ \\
$|V|$, $|U_1|$, $U$ (G19 cells) & Horizontal wind speed at the lowest level (surface-layer input); in the G19 port the local horizontal wind speed used as velocity scale. & m\,s$^{-1}$ \\
$V_{10}$ & 10-m wind speed (in the Rossby number of YSU). & m\,s$^{-1}$ \\
$\Theta$, $\theta$; $\Theta_v$, $\theta_v$; $\Theta_l$, $\theta_l$; $\theta_m$ & Potential temperature (mean, fluctuation); virtual; liquid-water; WRF's moist potential temperature $\theta_m = \theta(1 + 1.61 q_v)$. & K \\
$\theta_T$ & YSU's thermal excess added to $\theta_v$ of the first level in the bulk Richardson number. & K \\
$\Delta\theta_v$, $\Delta\theta$, $\Delta u$, $\Delta q$ & Jumps across the boundary-layer top (entrainment). & K; m\,s$^{-1}$; kg\,kg$^{-1}$ \\
$Q_v$, $q_v$; $Q_w$, $q_w$ & Water-vapour mixing ratio (mean, fluctuation); total water. & kg\,kg$^{-1}$ \\
$C$, $c$ & A generic transported scalar and its fluctuation (YSU flux formula). & as the scalar \\
$p$ & Pressure fluctuation. & Pa \\
$\rho$, $\rho_0$, $\rho_1$ & Air density; reference; at the lowest level. & kg\,m$^{-3}$ \\
$\avg{u_iu_j}$, $\tau_{ij}$; $\avg{u^2}$, $\avg{v^2}$, $\avg{w^2}$, $\avg{uv}$, $\avg{uw}$, $\avg{vw}$ & Reynolds (subgrid) stress tensor and its six components; $\tau_{ij}$ is the deviatoric form of Ito et al. & m$^2$\,s$^{-2}$ \\
$\avg{u_j\theta_v}$, $\avg{u\theta_v}$, $\avg{v\theta_v}$, $\avg{w\theta_v}$; $\avg{w\theta}$, $\avg{wq}$, $\avg{wc}$; $f_i$, $f_x$, $f_y$, $f_z$ & Turbulent heat-flux vector and its components; moisture and scalar fluxes; $f_i$ is the heat-flux vector of Ito et al. & K\,m\,s$^{-1}$; kg\,kg$^{-1}$\,m\,s$^{-1}$ \\
$\avg{\theta_v^2}$, $\avg{\theta^2}$ & Temperature variance. & K$^2$ \\
HFX, QFX & WRF surface sensible heat flux and moisture flux; $\avg{w\theta}_0 = \mathrm{HFX}/(\rho c_p)$. & W\,m$^{-2}$; kg\,m$^{-2}$\,s$^{-1}$ \\
$\tau_{xz}$ & Surface momentum flux (kinematic stress), $= u_*^2\,u_1/|U_1|$. & m$^2$\,s$^{-2}$ \\
\notgroup{Turbulence energy budget}
$e$ & Turbulence kinetic energy, $e = \tfrac12(\avg{u^2} + \avg{v^2} + \avg{w^2})$ (BouLac, SMS-3DTKE, LES, Ito). & m$^2$\,s$^{-2}$ \\
$\qsq$, $q$, $q_t$ & Twice the TKE (the prognostic variable of the Mellor--Yamada schemes; WRF calls it \texttt{qke}); $q = \sqrt{\qsq}$ is the turbulent velocity scale; $q_t = \sqrt{\qsq + q_{pe}}$ is WRF-MYNN's total turbulent velocity including the potential energy $q_{pe}$ (zero below closure 2.7). & m$^2$\,s$^{-2}$; m\,s$^{-1}$ \\
$E$, $R$ (Ito) & Resolved (grid-scale) kinetic energy of the deviations from the horizontal mean; subgrid fraction $R = e/(e + E)$. & m$^2$\,s$^{-2}$; -- \\
$P_s$, $P_b$, $\varepsilon$ & Shear production, buoyancy production, dissipation rate of TKE. & m$^2$\,s$^{-3}$ \\
$P_{HSP}$, hsp, \texttt{qHSP} (code) & Horizontal shear production added by the Goger ports. & m$^2$\,s$^{-3}$ \\
$P_{ke}$, pdk (code) & MYNN's production term of $\qsq/2$ at a level; rp $= \mathrm{pdk}_k + \mathrm{pdk}_{k+1}$ is what enters the $\qsq$ tendency. & m$^2$\,s$^{-3}$ \\
sh, bu (code) & BouLac's shear and buoyancy production per level. & m$^2$\,s$^{-3}$ \\
agg\_def2, $D_h^2$, $|D_h|$, $|S|_h$ & Horizontal deformation invariants: agg\_def2 $= D_h^2 = (\pd{U}{x})^2 + (\pd{V}{y})^2 + \tfrac12(\pd{U}{y} + \pd{V}{x})^2$ (Goger source); $|S|_h^2 = \tfrac14(D_{11} - D_{22})^2 + \overline{D_{12}}^{\,2}$ (WRF Smagorinsky, deviatoric). & s$^{-2}$; s$^{-1}$ \\
$D_{ij}$, $D_{11}, D_{12}, \ldots, D_{33}$ & Deformation tensor $D_{ij} = \pd{U_i}{x_j} + \pd{U_j}{x_i}$ ($D_{11} = 2\,\pd{U}{x}$). & s$^{-1}$ \\
$|S|$ & Strain-rate magnitude, $|S|^2 = 2(U_x^2 + V_y^2 + W_z^2) + (U_y + V_x)^2 + (U_z + W_x)^2 + (V_z + W_y)^2$. & s$^{-1}$ \\
$\Pi_{ij}$ & Pressure--strain correlation (return to isotropy). & m$^2$\,s$^{-3}$ \\
$\Gamma_\theta$, $\Gamma_v$ & NN09 level-3 countergradient term for heat; the corresponding term in MYNN's TKE production (dimensionless there, alongside $S_M G_M$). & K\,m$^{-1}$; -- \\
$\gamma_\theta$, $\gamma_q$, $\gamma_u$, $\gamma_v$, $\gamma_c$ & YSU / Shin--Hong countergradient amplitudes for heat, moisture, momentum, a generic scalar; the gradient correction actually used is $\gamma/h$. & K; kg\,kg$^{-1}$; m\,s$^{-1}$ \\
$\gamma$ (BouLac) & Troen--Mahrt countergradient term, $10\,\avg{w\theta}_0/(w_* h)$. & K\,m$^{-1}$ \\
$a_u(1)$, $b_\theta(1)$, $b_q(1)$ & BouLac's implicit drag coefficient and explicit flux sources of the lowest level. & s$^{-1}$; K\,s$^{-1}$; kg\,kg$^{-1}$\,s$^{-1}$ \\
\notgroup{Closure: eddy coefficients, stability functions, constants}
$K$, $K_M$, $K_H$; $K_m$, $K_h$, $K_c$, $K_q$ & Eddy viscosity (momentum) and eddy diffusivity (heat); $K = L q S$ in the Mellor--Yamada schemes; lower-case pairs in the YSU and Shin--Hong convention (see the index list). & m$^2$\,s$^{-1}$ \\
$K_{mh}$, $K_{mv}$, $K_{hh}$, $K_{hv}$; $K^S_{mh}$ & WRF's horizontal and vertical coefficients for momentum and heat; the Smagorinsky part of $K_{mh}$. & m$^2$\,s$^{-1}$ \\
$K_q$, $K_e$, $K_{q_{ke}}$ & Diffusivity of TKE ($K_q = S_q l q$ in the Mellor--Yamada schemes, $3K_m$ in WRF-MYNN, $2K$ in the Deardorff-type schemes). & m$^2$\,s$^{-1}$ \\
$K_0$ & YSU's local free-atmosphere coefficient before the Richardson-number factor, $K_0 = \ell^2|S|$. & m$^2$\,s$^{-1}$ \\
$S_M$, $S_H$, $S_q$, $S_\theta$ & Dimensionless stability functions (diffusivity ratios) for momentum, heat, TKE, temperature variance. & -- \\
$G_M$, $G_H$ & Non-dimensional vertical shear and buoyancy, $G_M = (L^2/q^2)[(\pd{U}{z})^2 + (\pd{V}{z})^2]$, $G_H = -(L^2/q^2)\,\beta g\,\pd{\Theta_v}{z}$. & -- \\
$E_1, \ldots, E_5$ & Auxiliary functions of $G_M$, $G_H$ in the level-2.5 stability functions. & -- \\
$F_1$, $F_2$, $f_1, \ldots, f_4$, $c_m$, $c_{h2}$ & Combinations of the closure constants in the level-2 functions (NN09 form; the fork's \texttt{fac1}--\texttt{fac4}, \texttt{cm}, \texttt{ch2}). & -- \\
$\alpha_c$, $q_{div}$ & Helfand--Labraga disequilibrium factor, $q/q_{eq}$ where $q < q_{eq}$, else 1. & -- \\
$A_1, A_2, B_1, B_2, C_1, \ldots, C_5$; $a_1, a_2, b_2, c_1$ & Mellor--Yamada closure constants (Rotta and dissipation time scales, isotropisation of production, pressure--scalar terms); lower case: the same constants in the fork's code. & -- \\
$\gamma_1$, $\gamma_2$ & NN09 combinations, $\gamma_1 = \tfrac13 - 2A_1/B_1$, $\gamma_2 = \tfrac{B_2}{B_1}(1 - C_3) + \tfrac{2A_1}{B_1}(3 - 2C_2)$; $\Rfc = \gamma_1/(\gamma_1 + \gamma_2)$. & -- \\
$a_{2fac}$, CK$_{mod}$ & Canuto/Kitamura factor on $A_2$, $1/(1 + \mathrm{CK}_{mod}\max(\Ri, 0))$, $\mathrm{CK}_{mod} = 1$. & -- \\
$\Pr$, $\Pr_0$, $\Pr_{lim}$ & Turbulent Prandtl number $K_M/K_H$; YSU's surface-layer value; MYNN's Richardson-dependent cap. & -- \\
$c_k$, $c_s$; $C_k$, $C_\varepsilon$; $c_{\varepsilon1}$, $c_{\varepsilon2}$ & Deardorff and Smagorinsky constants ($0.15$, $0.25$); BouLac's constants ($0.4$, $1/1.4$); WRF's dissipation constants ($C_\varepsilon = c_{\varepsilon1} + c_{\varepsilon2}\,\ell/\Delta$). & -- \\
$\alpha_1, \ldots, \alpha_6$, $c_{ns}$ & MYNN length-scale constants (turbulent-layer, buoyancy and stable-surface-layer coefficients). & -- \\
$\alpha$ & Three uses: the coefficient 0.10 of the Blackadar asymptotic length (Mellor--Yamada 1974, Juliano et al.); the wind-direction angle in the fork's surface stress ($\avg{uw} = -u_*^2\cos\alpha$); SMS-3DTKE's slope taper on the Smagorinsky coefficient. & -- \\
$b$, $a$; $\phi_{fac}$; $s_m$ & YSU countergradient constants ($a = b = 6.8$); the coefficient 8 in $w_s$; the Noh et al.\ momentum coefficient (15.9 in YSU and Shin--Hong, 10.9 in SMS-3DTKE). & -- \\
$c$ (Goger) & Smagorinsky constant of the horizontal shear production, 0.2. & -- \\
$\mathsf{A}$, $\mathbf{x}$, $\mathbf{b}$, $b_1$; $\mathsf{I}$, $\mathsf{S}$, $\tau^{-1}$ & The $10 \times 10$ linear system of the full closure: matrix, vector of the ten second moments, right-hand side and its first entry; the matrix is the identity times the inverse relaxation time plus the strain block. & mixed \\
$\mathcal{H}$ & Shin--Hong's nonlocal term written as a gradient correction, $\mathcal{H} = m_f/K_h$. & K\,m$^{-1}$ \\
$m_f$, $F_{nl}$, $\varepsilon_f$ & Shin--Hong's prescribed nonlocal heat-flux profile, its amplitude and the entrainment factor. & K\,m\,s$^{-1}$; K\,m\,s$^{-1}$; -- \\
\notgroup{Length and velocity scales}
$L$, $l$, $\ell$ & Master (mixing) length scale: the size of the energy-containing eddies. $L$ in NN09 and Ito et al., $l$ in the fork and BouLac, $\ell$ in WRF's diffusion code and YSU's local scheme. In $z/L$ and $L_{\mathrm{MO}}$, $L$ is the Obukhov length. & m \\
$L_S$, $L_T$, $L_B$; $l_S$, $l_T$, $l_B$, $l_F$ & MYNN blend: surface-layer, turbulent-layer, buoyancy and free-atmosphere lengths. & m \\
$l_0$, $\tilde q$ & Blackadar's asymptotic length ($q$-weighted mean height of the turbulence); the turbulent excess $\tilde q = \max(q - q_{min}, 0)$ that weights it in the fork. & m; m\,s$^{-1}$ \\
$l_{up}$, $l_{down}$, $l_k$, $l_\varepsilon$ & BouLac parcel displacements, mixing length $\min(l_{up}, l_{down})$, dissipation length $\sqrt{l_{up}l_{down}}$. & m \\
$l_M$, $l_{meso}$, $l_{LES}$, $\ell_v^{LES}$, $\ell_h$, $\ell_v$, $\ell_{diss}$, $\ell_{scale}$, $\Delta_{iso}$ & SMS-3DTKE and Deardorff lengths: mesoscale, LES-limit, horizontal, vertical, dissipation; WRF's \texttt{calc\_l\_scale}; isotropic grid length. & m \\
$l_{master}$, $l_{dissip}$, $l_{use}$, $l_{BouLac}$ & The fork's master scale, its dissipation copy, the strain-limited value entering the solve, the BouLac alternative. & m \\
$l_h$, $l_G$, $l_{eff}$ & Horizontal length of the Goger ports: $0.2\,\dx$ (G18); $l_G^2 = U^2 T_{L,u}T_{L,v}$ and its blended, capped form $l_{eff}$ (G19). & m \\
$L_x$, $L_\varepsilon$, $L_\theta$ & Ito et al.: horizontal length ($0.1\,h\,F$), dissipation length ($B_1 L$), temperature-variance dissipation length ($B_2 L$). & m \\
$\lambda$ & YSU's asymptotic length in the free atmosphere, 30--300~m. & m \\
$u_*$, $w_*$, $\theta_*$ & Friction velocity; convective velocity $w_* = (\beta g\,\avg{w\theta_v}_0\, z_i)^{1/3}$; surface temperature scale. & m\,s$^{-1}$; K \\
$w_s$, $w_{s0}$, $w_{s2}$ & YSU mixed-layer velocity scale $(u_*^3 + 8\kappa w_*^3\zeta)^{1/3}$; its value at $z = h/2$ (used for the countergradient terms); the top-down (cloud-top) branch. & m\,s$^{-1}$ \\
$w_m$ & Combined velocity scale $w_m^3 = w_*^3 + 5u_*^3$ (entrainment). & m\,s$^{-1}$ \\
$w_e$ & Entrainment velocity at the boundary-layer top. & m\,s$^{-1}$ \\
$q_c$ & NN09 convective velocity scale $(\beta g\,\avg{w\theta_v}_0\, L_T)^{1/3}$. & m\,s$^{-1}$ \\
$\sigma_u$, $\sigma_v$, $\sigma_w$ & Standard deviations of the velocity components (G19 port, from similarity profiles). & m\,s$^{-1}$ \\
$T_L$, $T_{L,u}$, $T_{L,v}$ & Lagrangian integral time scales, $0.15\,z_i/\sigma$. & s \\
$\tau$ & Relaxation time of the closure, $\propto l/q$ (row C7); in the fork's slope factor $\tau^2 = \tau_x^2 + \tau_y^2$ (dimensionless). & s; -- \\
$\delta$, $\delta_{sh}$ & Entrainment-zone depth (YSU, Shin--Hong). & m \\
\notgroup{Stability, boundary-layer depth, surface layer}
$g$, $g_j$, $\beta$ & Gravity; the gravity vector $(0, 0, -g)$ in the buoyancy terms of the stress equations; buoyancy coefficient $\beta = 1/T_{\mathrm{ref}}$, $T_{\mathrm{ref}} = 300$~K in every WRF scheme here. & m\,s$^{-2}$; K$^{-1}$ \\
$N$, $N^2$ & Brunt--V\"ais\"al\"a frequency, $N^2 = \beta g\,\pd{\Theta_v}{z}$. & s$^{-1}$; s$^{-2}$ \\
$\Ri$, $\Ri_b$ & Gradient Richardson number $N^2/[(\pd{U}{z})^2 + (\pd{V}{z})^2]$; bulk Richardson number from the surface (YSU, Shin--Hong, Goger 2019). & -- \\
$\Rf$, $\Rfc$, $R_{f1}$, $R_{f2}$ & Flux Richardson number $-P_b/P_s$; its critical value; two auxiliary values of the level-2 functions. & -- \\
$L_{\mathrm{MO}}$, $\zeta$ & Obukhov length; $\zeta = z/L_{\mathrm{MO}}$ in the MYNN length scales. In YSU $\zeta = (z - z_1)/(h - z_1)$, in the G19 port $\zeta = z/z_i$ (both dimensionless heights). & m; -- \\
$\phi_m$, $\phi_h$ & Monin--Obukhov stability functions for momentum and heat. & -- \\
$z_i$, $h$, $z_{i2}$, $h_1$, $h_2$, $z_{top}$ & Boundary-layer (mixed-layer) depth ($z_i$ in the Mellor--Yamada and SMS contexts, $h$ in YSU, Shin--Hong, BouLac, Ito); $z_{i2} = \max(z_i, 200~\mathrm{m})$ and the blend heights of MYNN's length scale; the model top. & m \\
$w_t$, $w_b$ & Blend weights: MYNN's BouLac blend above the PBL; the G19 port's blend to the Smagorinsky length above $\zeta = 0.8$. & -- \\
$\varphi_s$ & Shin--Hong: top of the surface layer as a fraction of $z_i$. & -- \\
$\mathrm{Ro}$, $f$, $z_0$ & Surface Rossby number $|V_{10}|/(f z_0)$; Coriolis parameter; roughness length. & --; s$^{-1}$; m \\
$C_t$, var\_sso, $\nabla^2 h$ & \texttt{topo\_wind} drag factor; WPS variance of the subgrid orography; terrain Laplacian as coded (five-point sum). & --; m$^2$; m \\
$v_{conv}$ & YSU's convective velocity that switches the \texttt{topo\_wind} correction off. & m\,s$^{-1}$ \\
$d_1$, $d_2$ & YSU entrainment-depth constants (0.02, 0.05). & -- \\
$c_p$, $c_p^*$, $c_{pm}$ & Specific heat of air at constant pressure; its moist forms $c_p(1 + 0.84 q_v)$ and $c_p(1 + 0.8 q_v)$. & J\,kg$^{-1}$\,K$^{-1}$ \\
$e_{min}$, $q_{min}$, TURB\_FLUX\_MIN & Numerical floors on TKE, on $q$ and on the fluxes. & m$^2$\,s$^{-2}$; m\,s$^{-1}$ \\
\notgroup{Grid, coordinates and scale awareness}
$\Delta$, $\dx$, $\Delta y$, $\Delta z$, $\overline{\Delta z}$ & Grid spacings; the mean of two adjacent layer thicknesses. & m \\
$\Delta_h$, $\Delta_{iso}$, $\Delta_s$, $\Delta_H$, $\Delta_{\mathrm{eff}}$ & Horizontal grid length $\sqrt{\dx\Delta y}$; isotropic length $(\dx\Delta y\Delta z)^{1/3}$ (Zhang et al.\ write $\Delta_s$); Ito's filter width; Shin--Hong's near-surface-inflated grid length $\Delta\max(0.2z_i/z, 1)$. & m \\
$\Delta^*$, $r$, $x$ & Grid length made dimensionless with the boundary-layer depth: $\Delta/h$ (Ito), $r = \Delta_h/z_i$ or $\max(2.5\Delta, 10)/\min(z_i, 3000)$ (SMS-3DTKE, MYNN, fork), $x = \Delta/(c_{sf}z_i)$ (Shin--Hong). & -- \\
$P(\cdot)$, $P_u$, $P_\theta$, $P_{nl}$, $P_{NL}$, $P_L$, $P_q$, $P_{TKE}$ & Partition functions between the LES limit (0) and the mesoscale limit (1). & -- \\
$F(\Delta^*)$ & Ito et al.'s length-scale factor, $P_{TKE}^{3/2}$. & -- \\
$\Psi_{bl}$, $\Psi_{shcu}$ & WRF-MYNN's scale-awareness tapers for the length scale and the mass flux. & -- \\
$c_{sf}$ & Shin--Hong stability factor on the nonlocal argument, between 1 and 2. & -- \\
$m$, $m_{tx}$, $m_{ux}$, $m_x$, $m_y$ & Map-scale factors (at mass and $u$ points). & -- \\
$\eta$ & WRF's terrain-following vertical coordinate; "along $\eta$" means along a model surface. & -- \\
$z_x$, $z_y$ & Slopes of the model surface, $(\pd{z}{x})_\eta$, $(\pd{z}{y})_\eta$. & -- \\
$\sigma$ & The fork's slope factor $\max(\tau, 1)$ dividing the horizontal tendencies; slope times grid aspect ratio. & -- \\
$\kappa$ & Von K\'arm\'an constant, 0.4. & -- \\
\notgroup{Scheme-specific remainders}
$a_{11}$, $a_{12}$, nlfrac, $f_{ent}$, $c_{ent}$ & Shin--Hong nonlocal-profile constants (1.0, $-1.15$, 0.7, 2, $-0.2$). & -- \\
$\delta_{ij}$ & Kronecker delta. & -- \\
$\pm$, $\mp$ in $\texttt{pbl3d\_opt} = \pm 1$ & The two approximate variants of the 3D closure: $-1$ applies vertical divergences only, $+1$ the full 3D divergences. & -- \\
\end{xltabular}
\normalsize

\section*{Legend}\label{sec:legend}
\begin{small}
\begin{tabularx}{\linewidth}{@{}>{\raggedright\arraybackslash}p{5.6cm} >{\raggedright\arraybackslash}X@{}}
\gK\ \gP\ \gD\ \gN & assumption kept \,/\, relaxed or partly dropped \,/\, dropped \,/\, not applicable to the scheme \\
\fglyph{\gD}\ {\setlength{\fboxsep}{1.6pt}\setlength{\fboxrule}{0.5pt}\fbox{\gP\ decisive cell}} & a \emph{decisive} cell: the assumption whose removal is the scheme's reason to exist, a cell that decides the questions of the first paper (horizontal production, $W$ gradients, length scale, stability cutoff), or a surprise \\
\colorbox{cShade}{\rule{0pt}{1.6ex}\ kept by every scheme\ } & shaded row: no scheme drops it --- the open frontier \\
\textcolor{cMYNN}{\textbf{MYNN as run}}\ \ \textcolor{cGXVIII}{\textbf{G18 port}}\ \ \textcolor{cGXIX}{\textbf{G19 port}} \newline \textcolor{cAPPROX}{\textbf{3D-APPROX}}\ \ \textcolor{cFULL}{\textbf{3D-FULL}} & the versions run in this project, in the colours of the project's diagnostics; colour means scheme identity and nothing else \\
\rg{day}\ \rg{night}\ \rg{both} & regime in which the failure matters: convective valley atmosphere, stable night, or both \\
\ev{obs}\ \ev{inf}\ \ev{run} & evidence: observed, inferred, or measured in this project's runs \\
\lit{} & literature evidence on the cell, added after the reading of 2026-09-07 \\
\fork\ \ \open & the fork or the WRF port differs from the publication \,/\, open item, named as such \\
\kfull\ \khalf & row the PhD concept turns on (a research question, regime or decision metric) \,/\, row the project's findings or methods depend on --- Section~I.3 reads the concept against the table \\
\end{tabularx}
\end{small}
